The established identities \(F=\mathcal D\) and \(R^{\star}\le F\) can be complemented by the following genuine-prior, rank-dependent converse.  Let \(M^{\star}\in\mathfrak M(K)\) maximize \(\mathcal D\), and suppose first that \(F=\mathcal D>0\).  Put \(r_{\star}=\operatorname{rank}(M^{\star})\), so \(r_{\star}\ge1\).  Because \(M^{\star}\in\operatorname{conv}\{v\otimes v:v\in K\}\), there are \(m<\infty\), numbers \(\lambda_k\ge0\) with \(\sum_{k=1}^m\lambda_k=1\), and schedules \(v^{(k)}\in K\) such that
\[
 M^{\star}=\sum_{k=1}^m\lambda_k v^{(k)}\otimes v^{(k)}.
\]
Define \(B:\mathbb R^m\to E\) by
\[
 Bt=\sum_{k=1}^m\sqrt{\lambda_k}\,t_kv^{(k)}.
\]
Then \(BB^{\ast}=M^{\star}\).  The set \(K\) in \(\mathrm{def:bounded\mbox{-}fourier\mbox{-}class}\) is convex and centrally symmetric and contains zero; hence it is absolutely convex.  Therefore, whenever \(\lVert t\rVert_2\le1\), the Cauchy--Schwarz inequality gives
\[
 \sum_{k=1}^m\left|\sqrt{\lambda_k}t_k\right|
 \le
 \left(\sum_{k=1}^m\lambda_k\right)^{1/2}
 \left(\sum_{k=1}^m t_k^2\right)^{1/2}
 \le1,
\]
and absolute convexity gives \(Bt\in K\).  Let \(H=(\ker B)^{\perp}\), let \(B_0\) be the restriction of \(B\) to \(H\), and identify \(H\) isometrically with \(\mathbb R^{r_{\star}}\).  Since \(B\) vanishes on \(\ker B\),
\[
 B_0B_0^{\ast}=BB^{\ast}=M^{\star},
 \qquad
 B_0\{t\in H:\lVert t\rVert_2\le1\}\subseteq K.
\]
Let \(U\) have normalized rotation-invariant probability law on the unit sphere of \(H\), and set \(\mathbf V=B_0U\).  This defines a countably additive symmetric probability law \(\nu_{\star}\in\mathcal P(K)\).  Rotation invariance gives
\[
 \mathbb E[U\otimes U]=r_{\star}^{-1}I_H,
 \qquad
 \mathbb E[\mathbf V\otimes\mathbf V]=r_{\star}^{-1}M^{\star}.
\]
Fix \(z\in\Omega_n\), and write \(A_z=N_zB_0\), \(b=B_0^{\ast}\ell\), and \(P_z=A_z^{\dagger}A_z\), the orthogonal projection of \(H\) onto \(\operatorname{range}(A_z^{\ast})\).  The value \(A_zU\) determines \(P_zU\).  Conditional on \(A_zU\), rotation invariance under the reflection that fixes \(\operatorname{range}(A_z^{\ast})\) and negates \(\ker A_z\) gives
\[
 \mathbb E[U\mid A_zU]=P_zU.
\]
Consequently the posterior mean of the target is linear in the observation:
\[
 \mathbb E[\theta(\mathbf V)\mid N_z\mathbf V]
 =\mathbb E[\langle b,U\rangle_H\mid A_zU]
 =\langle b,P_zU\rangle_H.
\]
Using \(\mathbb E[U\otimes U]=r_{\star}^{-1}I_H\),
\[
 \mathbb E_{\nu_{\star}}\!\left[
   \operatorname{Var}_{\nu_{\star}}(\theta(\mathbf V)\mid N_z\mathbf V)
 \right]
 =\frac{1}{r_{\star}}\lVert(I_H-P_z)b\rVert_2^2.
\]
On the other hand,
\[
 \langle\ell,M^{\star}\ell\rangle_E=\lVert b\rVert_2^2,
 \quad
 N_zM^{\star}\ell=A_zb,
 \quad
 N_zM^{\star}N_z^{\ast}=A_zA_z^{\ast},
\]
and the Moore--Penrose identity
\(A_z^{\ast}(A_zA_z^{\ast})^{\dagger}A_z=P_z\) yields
\[
 \psi_z(M^{\star})
 =\lVert b\rVert_2^2-\lVert P_zb\rVert_2^2
 =\lVert(I_H-P_z)b\rVert_2^2.
\]
Summing with the source masses specified by \(\mathrm{ass:source\mbox{-}product\mbox{-}design}\), whose positivity follows from \(\mathrm{ass:source\mbox{-}overlap}\), and applying \(\mathrm{def:prior\mbox{-}certificate}\), gives
\[
 \mathcal B(\nu_{\star})
 =\frac{1}{r_{\star}}
   \sum_{z\in\Omega_n}p_{\pi}(z)\psi_z(M^{\star})
 =\frac{\mathcal D}{r_{\star}}
 =\frac{F}{r_{\star}}.
\]
The proved proposition \(\mathrm{prop:primal\mbox{-}dual\mbox{-}certificates}\) therefore implies
\[
 \frac{F}{r_{\star}}\le R^{\star}\le F.
\]
If \(F=0\), the already established inequalities give \(R^{\star}=0\), so the same conclusion holds after replacing \(r_{\star}\) by \(r_{\star}\vee1\).  As an exact smallest-instance check, take one unit, one binary local exposure, complete local support, source success probability \(p\), and target success probability \(q\).  The independent Rademacher prior on the two local potential outcomes has a rank-two second moment and posterior ambiguity
\[
 (1-p)q^2+p(1-q)^2,
\]
which equals both frontiers in that instance; the spherical-ellipsoid construction gives the valid, but nonsharp, half of this value.  Thus the construction is a genuine symmetric prior and its calculation is consistent with the exact binary case, but its general comparison factor is \(1/r_{\star}\), not a numerical constant independent of the instance.